\documentclass[11pt]{article}

% ---------------------------------------------------------------------------
% Packages
% ---------------------------------------------------------------------------
\usepackage[margin=1in]{geometry}
\usepackage{lmodern}
\usepackage{amsmath,amssymb,amsthm,mathtools}
\usepackage{booktabs}
\usepackage{enumitem}
\usepackage{xcolor}
\usepackage[colorlinks=true,linkcolor=blue!55!black,citecolor=blue!50!black,urlcolor=blue!55!black]{hyperref}
\usepackage[expansion=false]{microtype}

% ---------------------------------------------------------------------------
% Theorem environments
% ---------------------------------------------------------------------------
\theoremstyle{plain}
\newtheorem{theorem}{Theorem}[section]
\newtheorem{proposition}[theorem]{Proposition}
\newtheorem{lemma}[theorem]{Lemma}
\newtheorem{corollary}[theorem]{Corollary}

\theoremstyle{definition}
\newtheorem{definition}[theorem]{Definition}
\newtheorem{model}[theorem]{Model}

\theoremstyle{remark}
\newtheorem*{intuition}{Intuition}
\newtheorem*{assumptionsnote}{Assumptions}
\newtheorem*{remark}{Remark}
\newtheorem*{example}{Example}

% Standing assumptions get their own counter (A1, A2, ...)
\newlist{standing}{enumerate}{1}
\setlist[standing]{label=\textbf{(A\arabic*)},leftmargin=2.6em}

% ---------------------------------------------------------------------------
% Math macros
% ---------------------------------------------------------------------------
\newcommand{\R}{\mathbb{R}}
\newcommand{\N}{\mathbb{N}}
\newcommand{\E}{\mathbb{E}}
\newcommand{\Prob}{\mathbb{P}}
\newcommand{\one}{\mathbb{1}}
\newcommand{\dd}{\,\mathrm{d}}
\newcommand{\Hist}{\mathcal{H}}
\newcommand{\Lap}{\mathcal{L}}
\newcommand{\Four}{\mathcal{F}}
\newcommand{\conv}{\ast}
\newcommand{\sfont}[1]{\mathsf{#1}}  % observed/data quantities in sans-serif
\newcommand{\xical}{\xi}
\newcommand{\Xical}{\Xi}
\DeclareMathOperator{\diag}{diag}
\DeclareMathOperator{\rank}{rank}
\DeclareMathOperator*{\argmin}{arg\,min}
\DeclareMathOperator{\KL}{KL}
\newcommand{\code}[1]{\texttt{#1}}
\newcommand{\half}{\tfrac12}

\title{\textbf{Interval-censored Hawkes Processes via the\\ Mean Behavior Poisson Process}\\[4pt]
\large From Poisson and Hawkes to covariate-modulated excitation,\\ and the Volterra systems that arise}
\author{A self-contained mathematical treatment\\
\normalsize companion to the \code{hawkes\_calibration} package}
\date{}

\begin{document}
\maketitle

\begin{abstract}
\noindent
The Hawkes process is the standard model for self-exciting events, but its
likelihood requires the exact event times. In many applications we observe only
\emph{interval-censored} data: counts of events per time bucket. This document
develops, from first principles and with full rigour, the machinery that makes
Hawkes calibration possible in that setting --- the \emph{Mean Behavior Poisson
process} (MBPP) --- and then pushes it through three levels of model complexity:
the bare process, time-varying \emph{baseline} covariates, and finally
time-varying \emph{excitation} covariates. The last is the hard case: it turns
the governing equation from a convolution into a \emph{general Volterra integral
equation of the second kind}, costing us the Laplace/Fourier shortcuts but still
solvable. Each result states its precise hypotheses and is accompanied by an
intuition. Appendices give self-contained tutorials on Volterra equations, the
Laplace and Fourier transforms, and ODE/state-space methods --- everything the
main text relies on.
\end{abstract}

\tableofcontents
\bigskip

\section{Introduction: the problem, and how we solve it}
\label{sec:intro}

\subsection{The problem}

A \emph{point process} records the times $t_1<t_2<\cdots$ of events. The
\emph{Hawkes process} is the canonical model when events beget events --- an
earthquake makes aftershocks more likely, a tweet makes re-tweets more likely.
Its defining feature is a \emph{self-exciting} conditional intensity
\begin{equation}
\lambda(t)=s(t)+\sum_{t_i<t}\phi(t-t_i),
\label{eq:intro-hawkes}
\end{equation}
where $s\ge 0$ is the exogenous ``background'' rate that injects \emph{immigrant}
events and the kernel $\phi\ge 0$ describes how each past event lifts the future
rate, producing \emph{offspring}. To fit \eqref{eq:intro-hawkes} by maximum
likelihood one needs every event time $t_i$.

In practice we frequently \emph{do not} observe event times. We observe
\emph{interval-censored} data: the \emph{number} of events in each time bucket ---
hospital admissions per day, video views per day, transactions per hour --- but
not when each event occurred, whether for privacy or because that is simply how
the data are recorded. Two facts then block a direct fit:
\begin{enumerate}[leftmargin=2.2em]
\item the log-likelihood of \eqref{eq:intro-hawkes} is \emph{undefined} without
the times $t_i$; and
\item the Hawkes process lacks the \emph{independent-increments} property, so we
cannot fall back on a Poisson likelihood for the bucket counts either (its
intensity, hence the number of events in a bucket, depends on the counts in
earlier buckets).
\end{enumerate}

\subsection{How we solve it}

The escape is to fit a \emph{different but parameter-matched} process whose mean
behaviour mimics the Hawkes process while enjoying independent increments. This is
the \emph{Mean Behavior Poisson process} (MBPP): the inhomogeneous Poisson process
whose deterministic intensity equals the \emph{expected} Hawkes intensity,
\begin{equation}
\xi(t):=\E[\lambda(t)]=s(t)+\int_0^t \phi(t-u)\,\xi(u)\dd u .
\label{eq:intro-mbpp}
\end{equation}
Three observations make this work, and they organise the whole document.
\begin{itemize}[leftmargin=1.4em]
\item \textbf{Why the MBPP is Poisson.} A simple point process whose compensator
is \emph{deterministic} is automatically Poisson (Watanabe's theorem,
\S\ref{sec:poisson-hawkes}); replacing the random Hawkes intensity $\lambda$ by its
mean $\xi$ removes the feedback and manufactures independent increments. The bucket
counts of the MBPP are therefore independent Poisson variables with means
$\Xi(o_{i-1},o_i]=\int_{o_{i-1}}^{o_i}\xi$, and their negative log-likelihood ---
the \emph{interval-censored log-likelihood} --- is a tractable sum of Poisson
terms.
\item \textbf{Why we can compute $\xi$.} Equation~\eqref{eq:intro-mbpp} is a
\emph{Volterra integral equation of the second kind}. Because $\phi$ here depends
only on the elapsed time $t-u$, it is a \emph{convolution}, so the system is
linear and time-invariant and we can solve it in closed form (an impulse
response), in the Laplace/Fourier domain (a transfer function), or as a linear
ODE. Appendix~\ref{app:volterra} is a self-contained primer on Volterra equations;
Appendices~\ref{app:transforms} and \ref{app:ode} cover the transform and ODE
methods.
\item \textbf{Why the parameters transfer.} The MBPP and the Hawkes process are
built from the \emph{same} $s$ and $\phi$, so fitting the MBPP on the counts
recovers (approximations of) the Hawkes parameters.
\end{itemize}

\subsection{Roadmap, in increasing order of difficulty}

We deliberately build up the model one layer at a time.
\begin{description}[leftmargin=1.4em,style=nextline]
\item[\S\ref{sec:poisson-hawkes} Poisson and Hawkes.] The two ingredients:
the Poisson process (and exactly why a deterministic compensator makes a process
Poisson) and the Hawkes process (and exactly why it is \emph{not} Poisson), with
the cluster/branching picture that yields \eqref{eq:intro-mbpp} rigorously.
\item[\S\ref{sec:assumptions} Standing assumptions.] The precise regularity
conditions under which everything below holds, stated once.
\item[\S\ref{sec:no-covariates} Without covariates.] The univariate and
multivariate MBPP. Here $\phi$ is a fixed convolution kernel, the equation is LTI,
and everything is closed-form. This is the workhorse.
\item[\S\ref{sec:baseline-cov} Baseline covariates.] Let the immigrant rate
depend on covariates, $s(t)=\exp(\gamma_0+\gamma^\top X(t))$. The equation stays a
convolution (only the forcing changes), so the LTI machinery survives unchanged;
we simply fit the coefficients $\gamma$.
\item[\S\ref{sec:excitation-cov} Excitation covariates (the hard case).] Let the
\emph{excitation itself} depend on covariates, $\alpha(t)=\alpha_0\exp(\delta^\top
Z(t))$. Now the kernel depends on $t$ and $u$ \emph{separately}, the equation is a
\emph{general (non-convolution) Volterra equation}, and time-invariance ---
together with the Laplace/Fourier shortcuts --- is lost. We show the problem is
still well-posed and, for exponential kernels, reduces to a \emph{linear
time-varying} ODE that we can integrate.
\end{description}

\noindent
The short answer to ``can we put covariates in the excitation?'' is
\emph{yes} --- \S\ref{sec:excitation-cov} does it --- but it is genuinely harder,
and seeing exactly \emph{why} is the reason the Volterra viewpoint is worth
developing carefully.

\section{Poisson and Hawkes}
\label{sec:poisson-hawkes}

We need two building blocks: the Poisson process (which the MBPP will be) and the
Hawkes process (which we want to fit). The crucial contrast is whether the
\emph{compensator} is deterministic.

\subsection{Counting processes, intensity, compensator}

A simple temporal point process is described by its counting process
$N(t)=\#\{i:t_i\le t\}$, adapted to its internal history
$\Hist_t=\sigma\{N(s):s\le t\}$. By the Doob--Meyer decomposition there is a unique
predictable increasing \emph{compensator} $\Lambda$ with
\begin{equation}
M(t)=N(t)-\Lambda(t)\ \text{an }(\Hist_t)\text{-local martingale.}
\end{equation}
When $\Lambda$ is absolutely continuous, $\Lambda(t)=\int_0^t\lambda$, and the
density $\lambda$ is the \emph{conditional intensity}
$\lambda(t)=\lim_{h\downarrow0}\frac1h\,\E[N(t{+}h)-N(t)\mid\Hist_{t^-}]$.

\begin{intuition}
$\Lambda$ is the predictable ``running forecast'' of how many events should have
happened by time $t$; $\lambda$ is that forecast's instantaneous rate. Subtracting
the forecast from reality leaves pure surprise --- the martingale $M$ --- so
$\lambda$ captures exactly the part of the future the past lets us anticipate.
\end{intuition}

\subsection{The Poisson process and why deterministic $\Lambda$ matters}

\begin{definition}[Poisson process]
$N$ is an inhomogeneous Poisson process with mean measure $\Lambda$ if for disjoint
bounded sets $A_1,\dots,A_k$ the counts $N(A_i)$ are independent with
$N(A_i)\sim\mathrm{Poisson}(\Lambda(A_i))$.
\end{definition}

Two consequences are used constantly: a Poisson process has \emph{independent
increments}, and $\E[N(t)]=\Lambda(t)$. The deep fact we rely on is the converse.

\begin{theorem}[Watanabe characterisation]\label{thm:watanabe}
Let $N$ be a simple point process with continuous compensator $\Lambda$. Then $N$
is an inhomogeneous Poisson process with mean measure $\Lambda$ \emph{if and only
if $\Lambda$ is deterministic.}
\end{theorem}

\begin{intuition}
Dependence between increments can only come from the intensity \emph{reacting to
history}. If the intensity --- hence the compensator --- is fixed in advance, there
is no channel through which the past can speak to the future, so the process must
be Poisson. \textbf{This single fact is why the MBPP works}: replacing the random
Hawkes intensity $\lambda$ by its deterministic mean $\xi$ severs the feedback loop
and gives us back independent increments (and therefore a product likelihood,
\S\ref{sec:no-covariates}).
\end{intuition}

\subsection{The likelihood of a point process}

\begin{theorem}[Jacod formula]\label{thm:jacod}
A simple point process with $(\Prob,\Hist_t)$-intensity $\lambda$ has, on $[0,T]$,
the log-likelihood (relative to a unit-rate Poisson reference, up to a constant)
\begin{equation}
\ell(\theta)=\sum_{t_i\le T}\log\lambda(t_i;\theta)-\int_0^T\lambda(s;\theta)\dd s .
\label{eq:ppll}
\end{equation}
\end{theorem}

\begin{intuition}
The likelihood rewards forecasting high intensity exactly where events landed
($\sum\log\lambda(t_i)$) and penalises total forecast mass ($\int\lambda$): ``put
probability on the bumps, don't waste it elsewhere.'' It is the point-process
analogue of $\prod_i p(x_i)$, and it manifestly needs the times $t_i$ --- the
obstruction in the interval-censored setting.
\end{intuition}

\subsection{The Hawkes process and why it is \emph{not} Poisson}

\begin{model}[Hawkes process]\label{mod:hawkes}
The Hawkes process has stochastic intensity
\begin{equation}
\lambda(t)=s(t)+\int_0^{t^-}\!\phi(t-u)\dd N(u)
=s(t)+\sum_{t_i<t}\phi(t-t_i),
\label{eq:hawkes}
\end{equation}
with exogenous (immigrant) rate $s$ and triggering kernel $\phi$. The
\emph{exponential kernel} is $\phi(t)=\kappa\theta\,e^{-\theta t}\,\one[t>0]$, for
which $\int_0^\infty\phi=\kappa$ is the \emph{branching ratio} and $\theta$ the
decay rate.
\end{model}

Because $\lambda$ depends on the realised history, the compensator
$\Lambda(t)=\int_0^t\lambda$ is \emph{random}; by Theorem~\ref{thm:watanabe} the
Hawkes process is therefore \emph{not} Poisson and has dependent increments.

\subsection{The cluster representation and the mean-intensity equation}

\begin{theorem}[Hawkes--Oakes cluster representation]\label{thm:cluster}
The Hawkes process of \eqref{eq:hawkes} is distributionally equal to the following
branching construction: \emph{(i)} immigrants arrive as a Poisson process of
intensity $s$; \emph{(ii)} each event at time $u$ independently produces direct
offspring as a Poisson process of intensity $\phi(\cdot-u)$ on $(u,\infty)$;
\emph{(iii)} the process is the superposition of all immigrants and descendants.
\end{theorem}

\begin{intuition}
Self-excitation made literal: every event founds a ``family,'' and the whole
process is a \emph{forest of independent family trees} rooted at immigrants. The
mean number of direct children per event is $\int\phi=\kappa$, so a family is a
Galton--Watson tree with mean offspring $\kappa$ --- finite on average iff
$\kappa<1$ (\emph{subcritical}), with expected total size $1/(1-\kappa)$.
\end{intuition}

Averaging the cluster representation gives the equation at the heart of everything.

\begin{proposition}[Mean-intensity / MBPP equation]\label{prop:mbpp-eq}
Under the standing assumptions of \S\ref{sec:assumptions}, the expected intensity
$\xi(t)=\E[\lambda(t)]$ satisfies
\begin{equation}
\xi(t)=s(t)+\int_0^t\phi(t-u)\,\xi(u)\dd u .
\label{eq:mbpp}
\end{equation}
\end{proposition}
\begin{proof}
Take $\E[\cdot]$ in \eqref{eq:hawkes}. By Tonelli (non-negative integrands) the
expectation passes inside, and using the tower property
$\E[\dd N(u)]=\E[\E[\dd N(u)\mid\Hist_{u^-}]]=\E[\lambda(u)]\dd u=\xi(u)\dd u$.
\end{proof}

\begin{intuition}
Average the \emph{stochastic} feedback over all realisations and it becomes
\emph{deterministic} feedback: ``on average, today's rate is the baseline plus the
kernel-weighted echo of the average past.'' Equation~\eqref{eq:mbpp} is the object
we will solve again and again. It is a \emph{Volterra integral equation of the
second kind} (Appendix~\ref{app:volterra}); when --- as here --- the kernel depends
only on $t-u$, it is moreover a \emph{convolution}, which is what makes the next
section entirely closed-form.
\end{intuition}

\begin{remark}[Why interval censoring is hard for Hawkes directly]
Decomposing a bucket count over offspring generations produces nested
\emph{Borel} convolutions with no finite closed form at finite $T$, and Monte
Carlo over the generation structure scales poorly. This is the formal sense in
which ``the Hawkes process cannot be fit interval-censored'' --- and the reason we
detour through the MBPP.
\end{remark}

\section{Standing assumptions for our model}
\label{sec:assumptions}

We collect, once, the precise regularity conditions under which every result
below holds. They are mild and satisfied by every model in the companion package
(exponential kernels; power-law kernels with positive offset; non-negative,
locally integrable baselines). Each later result states any \emph{extra}
hypothesis explicitly.

\begin{standing}
\item \textbf{Simplicity and non-explosion.} $N$ is a simple temporal point
process on $[0,\infty)$: no two events coincide ($\Prob$-a.s.), and $N(t)<\infty$
a.s.\ for every finite $t$.
\item \textbf{Stochastic intensity.} $N$ admits an $(\Hist_t)$-predictable
conditional intensity $\lambda\ge 0$ with respect to its internal history
$\Hist_t=\sigma\{N(s):s\le t\}$; equivalently the compensator
$\Lambda(t)=\int_0^t\lambda$ is absolutely continuous, with $\int_0^T\lambda<\infty$
a.s.
\item \textbf{Kernel.} $\phi:\R\to[0,\infty)$ is measurable, causal ($\phi(u)=0$
for $u<0$), and locally bounded, with branching ratio
$n^{*}:=\int_0^\infty\phi=\kappa<\infty$. \emph{Subcritical} means $\kappa<1$.
\item \textbf{Baseline.} $s:[0,\infty)\to[0,\infty)$ is measurable and locally
integrable. Where a Dirac-comb (multi-impulse) baseline appears it is a
non-negative measure and convolutions with it are read distributionally.
\item \textbf{Finite mean intensity.} Wherever $\E[\lambda]$ is used,
$\sup_{t\le T}\E[\lambda(t)]<\infty$ --- automatic under (A3)-subcritical with
locally integrable $s$.
\item \textbf{Smooth, identifiable parameterisation.} For estimation, $\theta$
ranges over an open $\Theta\subseteq\R^p$; the map $\theta\mapsto\lambda(t;\theta)$
is twice continuously differentiable with $\lambda(t;\theta)>0$ on the data, and
the model is identifiable on $\Theta$.
\item \textbf{Multivariate.} In dimension $M$: $N=(N_1,\dots,N_M)$, intensity
$\lambda\in\R_+^M$, kernel \emph{matrix} $\Phi=(\phi_{mj})$, and \emph{branching
matrix} $G=(\int\phi_{mj})$. \emph{Stable} means the spectral radius
$\rho(G)<1$.
\end{standing}

\begin{intuition}
(A1)--(A2) say the process is well-behaved enough to have an instantaneous rate;
(A3)--(A4) say the two ingredients $\phi$ and $s$ are honest non-negative
functions with finite total excitation; (A5) keeps means finite; (A6) makes
maximum likelihood meaningful; (A7) lifts everything to $M$ interacting
components, where the single number $\kappa<1$ is replaced by the matrix condition
$\rho(G)<1$. Whenever a later theorem needs more --- boundedness of $\phi$,
subcriticality, existence of a transform --- it is flagged at that result.
\end{intuition}

\begin{remark}[Two parameterisations of the exponential kernel]
The package's event-time code writes $\phi(t)=\alpha e^{-\beta t}$, while the MBPP
theory uses $\phi(t)=\kappa\theta e^{-\theta t}$. They coincide with
$\alpha=\kappa\theta,\ \beta=\theta$, i.e.\ $\kappa=\alpha/\beta$ \emph{is} the
branching ratio. We use $(\kappa,\theta)$ throughout the theory.
\end{remark}

\section{The MBPP without covariates}
\label{sec:no-covariates}

Here $\phi$ is a fixed convolution kernel. The MBPP equation~\eqref{eq:mbpp} is a
\emph{convolution} Volterra equation of the second kind, hence linear and
time-invariant (LTI), and everything is closed-form. This is the workhorse the
rest of the paper perturbs.

\subsection{Existence, uniqueness, and the resolvent}

Write the Volterra operator $(\mathcal V\xi)(t)=\int_0^t\phi(t-u)\xi(u)\dd u$, so
\eqref{eq:mbpp} reads $(\mathcal I-\mathcal V)\xi=s$.

\begin{theorem}[Well-posedness; see Appendix~\ref{app:volterra}]\label{thm:exist}
If $\phi$ is bounded on $[0,T]$ (true for the exponential kernel and for power-law
with positive offset) and $s\in C[0,T]$, then \eqref{eq:mbpp} has a unique
continuous solution on $[0,T]$, given by the Neumann series
$\xi=\sum_{n\ge0}\mathcal V^n s$.
\end{theorem}

\begin{intuition}
On a finite horizon, causal feedback can only compound so fast: the $n$-th echo
$\mathcal V^n$ carries a $t^n/n!$ factor that crushes it, so the loop
$\mathcal I-\mathcal V$ is \emph{always} invertible there, whatever the gain
$\kappa$. ``No runaway in finite time.'' Subcriticality ($\kappa<1$) is needed only
to extend this to $t=\infty$.
\end{intuition}

Separating the $n=0$ term, $\xi=s+h\conv s$ with the \emph{resolvent kernel}
$h=\sum_{n\ge1}\phi^{\conv n}$ ($\phi^{\conv n}$ the $n$-fold convolution). The
companion impulse response is $E=\delta+h$, so $\xi=E\conv s$
(Appendix~\ref{app:volterra}, \ref{app:transforms}).

\begin{proposition}[Global resolvent]\label{prop:resolvent}
If $\kappa=\int_0^\infty\phi<1$ then $h\in L^1(0,\infty)$ with
$\|h\|_1=\kappa/(1-\kappa)$, and $\xi=s+h\conv s$ holds on $[0,\infty)$.
\end{proposition}

\begin{intuition}
The resolvent sums the echo over \emph{all} offspring generations; it is globally
integrable exactly when each generation shrinks ($\kappa<1$), and its total mass
$\kappa/(1-\kappa)$ is the expected number of offspring per event. The resolvent
\emph{is} the renewal density of the branching tree.
\end{intuition}

\subsection{The exponential kernel closes everything}

For $\phi(t)=\kappa\theta e^{-\theta t}$ one finds (Appendix~\ref{app:ode})
\begin{equation}
h(t)=\sum_{n\ge1}\phi^{\conv n}(t)=\kappa\theta\,e^{(\kappa-1)\theta t},\qquad t>0,
\label{eq:h-exp}
\end{equation}
a single exponential (the series telescopes because one timescale stays one
timescale). The compensator obeys the \emph{same} equation with $s$ replaced by its
integral $\sfont S(t)=\int_0^t s$:
\begin{equation}
\Xi(t)=\sfont S(t)+\int_0^t\phi(t-u)\,\Xi(u)\dd u=\sfont S+h\conv\sfont S,
\qquad \Xi(x,y]=\Xi(y)-\Xi(x).
\end{equation}

\begin{example}[A single immigrant]
For $s=\delta(\cdot-a)$ (one immigrant at $a$), with $c:=(\kappa-1)\theta<0$,
\begin{equation*}
\xi(t)=\delta(t-a)+\kappa\theta\,e^{c(t-a)}\one[t>a],\qquad
\Xi^{\mathrm{endo}}(t)=\frac{\kappa}{1-\kappa}\bigl(1-e^{c(t-a)}\bigr)\one[t>a].
\end{equation*}
As $t\to\infty$, $\Xi^{\mathrm{endo}}\to\kappa/(1-\kappa)$: the mean offspring count
of one immigrant. The lone event leaves a \emph{decaying wake} of expected
offspring whose total mass is the branching sum.
\end{example}

\subsection{Multivariate MBPP}

With $M$ interacting components, $\xi\in\R^M$ solves the matrix convolution Volterra
equation
\begin{equation}
\xi(t)=s(t)+\int_0^t\Phi(t-u)\,\xi(u)\dd u,\qquad \Phi=(\phi_{mj}),
\label{eq:mbpp-mv}
\end{equation}
where $\phi_{mj}$ is how component $j$ excites component $m$. In the
transform domain (Appendix~\ref{app:transforms}) this is
$\hat\xi(\omega)=(I-\hat\Phi(\omega))^{-1}\hat s(\omega)$.

\begin{proposition}[Stationarity and spectral radius]\label{prop:stationary}
Under a constant baseline $s\equiv\bar\mu$ and stability $\rho(G)<1$ (where
$G_{mj}=\int\phi_{mj}=\hat\Phi(0)_{mj}$ is the branching matrix), the stationary
mean intensity is
\begin{equation}
\bar\xi=(I-G)^{-1}\bar\mu=\sum_{n\ge0}G^n\bar\mu .
\end{equation}
\end{proposition}
\begin{proof}
Stationarity in \eqref{eq:mbpp-mv} gives $\bar\xi=\bar\mu+G\bar\xi$; the Neumann
series converges iff $\rho(G)<1$.
\end{proof}

\begin{intuition}
The stationary intensity is the baseline amplified by the network's \emph{total}
feedback $(I-G)^{-1}=I+G+G^2+\cdots$: direct excitation, plus two-hop, plus
three-hop, \dots. The condition $\rho(G)<1$ is exactly what makes those multi-hop
echoes sum to something finite.
\end{intuition}

For exponential kernels the matrix equation is realised, with no transform, as a
finite \emph{linear ODE} in $M^2$ states $y_{mj}=\bigl(a_{mj}e^{-b_{mj}\cdot}\bigr)\conv\xi$
(Appendix~\ref{app:ode}):
\begin{equation}
y_{mj}'=a_{mj}\,\xi_j-b_{mj}\,y_{mj},\qquad \xi_m=s_m+\sum_j y_{mj},
\label{eq:statespace}
\end{equation}
with $a_{mj}=\kappa_{mj}\theta_{mj}$, $b_{mj}=\theta_{mj}$. This is the workhorse
solver \code{solve\_mbpp\_ode\_multivariate}; the numerical check
$\bar\xi=(I-G)^{-1}\bar\mu$ holds to machine precision.

\subsection{Fitting on interval-censored counts (IC-LL)}
\label{sub:icll}

Because the MBPP is Poisson (Theorem~\ref{thm:watanabe}), the bucket counts over an
observation grid $0=o_0<\cdots<o_m=T$ are \emph{independent}
$\mathrm{Poisson}(\Xi(o_{i-1},o_i])$. The negative log-likelihood is therefore a
simple sum.

\begin{proposition}[Interval-censored log-likelihood]\label{prop:icll}
For observed counts $\sfont C(o_{i-1},o_i]$,
\begin{equation}
\mathcal L_{\mathrm{IC}}(\theta)=\sum_{i=1}^m\Xi(o_{i-1},o_i;\theta]
-\sum_{i=1}^m \sfont C(o_{i-1},o_i]\,\log\Xi(o_{i-1},o_i;\theta],
\label{eq:icll}
\end{equation}
up to a constant in $\theta$. Minimising it over the kernel parameters recovers the
Hawkes parameters through the MBPP correspondence.
\end{proposition}

\begin{intuition}
Independent Poisson buckets $\Rightarrow$ the joint probability is a product
$\Rightarrow$ the log-likelihood is a sum of per-bucket Poisson terms in the
interval compensators $\Xi_i$. It compares a \emph{count} $\sfont C_i$ to an
\emph{expected count} $\Xi_i$ (both extensive), as it should. (Choosing a Gaussian
noise model instead of Poisson turns the loss into least squares; see
Appendix~\ref{app:estimation} for the Bregman/exponential-family picture and the
identifiability of $\kappa$ versus $\theta$.)
\end{intuition}

\begin{remark}[What is robustly estimable]
The branching ratio $\kappa$ (the gain) is well identified from counts; the
timescale $\theta$ is weakly identified, because a bucket integrates away the
within-bucket timing that carries $\theta$
(Appendix~\ref{app:estimation}). Treat $\kappa$ as trustworthy and report $\theta$
with wide intervals.
\end{remark}

\section{Adding covariates to the baseline}
\label{sec:baseline-cov}

The first and easiest place to let data drive the model is the \emph{baseline}
(immigrant) rate. Let observed covariates $X(t)\in\R^p$ enter through a log-linear
link,
\begin{equation}
s(t)=\exp\!\bigl(\gamma_0+\gamma^\top X(t)\bigr),
\label{eq:loglinear}
\end{equation}
which keeps $s\ge0$ for all $\gamma$.

\subsection{Why this stays easy}

The MBPP equation~\eqref{eq:mbpp} is \emph{linear in the forcing $s$} and its
kernel is unchanged: \eqref{eq:loglinear} only changes \emph{what} we feed in, not
the operator $\mathcal I-\mathcal V$. So the entire LTI solution of
\S\ref{sec:no-covariates} carries over verbatim, $\xi=E\conv s$ and
$\Xi=E\conv\sfont S$ with $\sfont S=\int_0^t s$.

\begin{proposition}[Piecewise-constant covariates remain closed-form]
\label{prop:pc-cov}
If $X(t)$ is piecewise-constant, then $s(t)=\exp(\gamma_0+\gamma^\top X(t))$ is
piecewise-constant, with rate $\exp(\gamma_0+\gamma^\top X_k)$ on the $k$-th
covariate interval. The MBPP intensity and compensator are then the same
closed-form piecewise-constant responses as in \S\ref{sec:no-covariates}; only the
per-interval \emph{rates} are reparameterised.
\end{proposition}

\begin{intuition}
A covariate that is constant on each bucket simply re-labels the height of the
baseline on that bucket. The self-excitation machinery does not know or care that
the height came from a regression rather than a free parameter --- so nothing about
the solution method changes, and the closed forms of \S\ref{sec:no-covariates}
apply unchanged.
\end{intuition}

\subsection{Fitting the covariate effects}

We now optimise \eqref{eq:icll} over $(\gamma_0,\gamma)$ \emph{jointly} with the
kernel parameters $(\kappa,\theta)$. The compensator $\Xi$ is assembled from the
log-linear baseline exactly as before, so the only change from
\S\ref{sub:icll} is a larger but still smooth parameter vector. This is the
interval-censored analogue of the event-time covariate model, implemented as
\code{fit\_mbpp\_ic\_covariates}.

\begin{remark}[Tabular data]
A time-bucketed table --- one row per observation interval, with covariate columns
$X$ and an event-count column $\sfont C$ --- maps directly onto the fit: the row
endpoints are the observation grid, $\sfont C$ is the count vector, and the
covariate columns are $X$. Categorical covariates are one-hot encoded first.
\end{remark}

\subsection{What is and is not identified}

Empirically and theoretically (Appendix~\ref{app:estimation}), the covariate
\emph{effects} $\gamma$ are recovered tightly, because they are read off from how
counts change \emph{across} buckets as $X$ varies; the absolute split between the
baseline level $\gamma_0$ and the self-excitation $\kappa$ is more weakly
identified, sharing the $\kappa$--$\theta$ ridge. A covariate that does not vary
across buckets carries no information about its coefficient: identification of
$\gamma$ requires covariate \emph{contrast}.

\begin{intuition}
The model can explain a given level of activity either as ``a high baseline'' or
``a modest baseline strongly amplified.'' Counts alone struggle to separate the
two. But \emph{differences} in activity that track \emph{differences} in a
covariate pin that covariate's effect cleanly --- which is why $\gamma$ is the
trustworthy output here.
\end{intuition}

\medskip
\noindent
Crucially, in \eqref{eq:loglinear} the kernel $\phi$ is still a fixed convolution
kernel: covariates have entered only the \emph{forcing}. The next section asks what
happens when they enter the \emph{kernel} --- and that is where the convolution
structure, and with it the LTI shortcuts, finally break.

\section{Excitation covariates: the Volterra system that arises}
\label{sec:excitation-cov}

We now let covariates modulate the \emph{excitation itself} --- the strength with
which one event triggers others. This is the genuinely harder case, and the whole
point of developing the Volterra viewpoint: the governing equation stops being a
convolution.

\subsection{The model}

\begin{model}[Covariate-modulated excitation]\label{mod:excitation}
Let observed covariates $Z(t)\in\R^q$ scale the triggering strength through a
log-linear link,
\begin{equation}
\alpha_{mj}(t)=\alpha^{0}_{mj}\,\exp\!\bigl(\delta_{mj}^{\top}Z(t)\bigr),
\qquad
\phi_{mj}(t,u)=\alpha_{mj}(t)\,g_{mj}(t-u),
\label{eq:excit-kernel}
\end{equation}
where $g_{mj}\ge0$ is a fixed shape (e.g.\ $g_{mj}(\tau)=\theta_{mj}e^{-\theta_{mj}\tau}$)
and $\alpha^0_{mj}\ge0$ is the covariate-free base strength.
\end{model}

The modulation may be read at the \emph{receiving} time $t$ (as written) or at the
\emph{parent} time $u$ (replace $Z(t)$ by $Z(u)$); both are handled below and the
implementation uses the former.

\subsection{The mean-intensity equation is now a general Volterra equation}

Averaging the Hawkes intensity exactly as in Proposition~\ref{prop:mbpp-eq}, but
with the time-dependent kernel \eqref{eq:excit-kernel}, gives

\begin{proposition}[Non-convolution MBPP equation]\label{prop:nonconv}
Under (A1)--(A7), the expected intensity solves the system of linear Volterra
equations of the second kind
\begin{equation}
\xi_m(t)=s_m(t)+\sum_{j=1}^M\int_0^t K_{mj}(t,u)\,\xi_j(u)\dd u,
\qquad
K_{mj}(t,u)=\alpha^0_{mj}\,e^{\delta_{mj}^\top Z(t)}\,g_{mj}(t-u).
\label{eq:mbpp-nonconv}
\end{equation}
The kernel $K_{mj}(t,u)$ depends on $t$ and $u$ \emph{separately}, not only on
$t-u$.
\end{proposition}

\begin{intuition}
In \S\ref{sec:no-covariates}--\ref{sec:baseline-cov} the kernel was a function of
the elapsed time $t-u$ alone --- the system reacted to a stimulus the same way
regardless of \emph{when} it arrived. Now the response strength rides on $Z(t)$, so
\emph{when} matters: the same past event echoes more loudly in a high-$Z$ period
than in a low-$Z$ period. Mathematically the kernel is no longer a convolution, and
\textbf{this is exactly what costs us the LTI toolkit}.
\end{intuition}

\subsection{What survives, and what is lost}

\begin{description}[leftmargin=1.4em,style=nextline]
\item[Well-posedness survives.] Equation~\eqref{eq:mbpp-nonconv} is a standard
linear Volterra system. By the general theory (Appendix~\ref{app:volterra}, which
does \emph{not} use convolution structure), for any locally bounded kernel and
continuous $s$ it has a unique solution on every finite horizon, given by the
Neumann series $\xi=\sum_{n\ge0}\mathcal K^n s$ with
$(\mathcal K\xi)_m(t)=\sum_j\int_0^t K_{mj}(t,u)\xi_j(u)\dd u$.
\item[Time-invariance is lost.] Because $K(t,u)\ne K(t-u)$, the convolution theorem
no longer applies: there is \emph{no} Laplace/Fourier solution, \emph{no} constant
transfer function $R(\omega)$, and the spectral (Fourier-neural-operator) shortcut
of the convolution case is no longer exact. The resolvent is now a genuine
two-variable kernel $R(t,u)$ solving $R=K+K\!\circ\!R$ (Appendix~\ref{app:volterra}),
not a single function of $t-u$.
\end{description}

\subsection{Exponential shape $\Rightarrow$ a finite linear time-varying ODE}

All is not lost computationally. The key is that an exponential $g$ makes the
kernel \emph{separable} (degenerate): with $g_{mj}(\tau)=e^{-b_{mj}\tau}$,
\begin{equation*}
K_{mj}(t,u)=\underbrace{\alpha^0_{mj}\,e^{\delta_{mj}^\top Z(t)}\,e^{-b_{mj}t}}_{a_{mj}(t)}\;\underbrace{e^{\,b_{mj}u}}_{\text{(}u\text{ only)}},
\end{equation*}
a rank-one product of a $t$-function and a $u$-function. A separable Volterra
kernel collapses the integral equation onto finitely many ODEs
(Appendix~\ref{app:volterra}). Concretely, introducing the filtered histories
$y_{mj}(t)=\int_0^t e^{-b_{mj}(t-u)}\xi_j(u)\dd u$:

\begin{theorem}[LTV state-space reduction]\label{thm:ltv}
For exponential shapes, the non-convolution MBPP equation~\eqref{eq:mbpp-nonconv}
(receiving-time modulation) is equivalent to the finite \emph{linear time-varying}
ODE system
\begin{equation}
y_{mj}'(t)=\xi_j(t)-b_{mj}\,y_{mj}(t),
\qquad
\xi_m(t)=s_m(t)+\sum_{j} \alpha^0_{mj}\,e^{\delta_{mj}^\top Z(t)}\,y_{mj}(t),
\label{eq:ltv}
\end{equation}
with $y_{mj}(0)=0$. The coefficients are constant except for the modulation
$e^{\delta^\top Z(t)}$, which makes the system time-\emph{varying} rather than
time-invariant.
\end{theorem}

\begin{intuition}
The history of each component is still summarised by a finite ``running, decaying
memory'' $y_{mj}$ --- exactly the state of the convolution case
\eqref{eq:statespace}. The \emph{only} change is that the gain with which that
memory feeds back, $\alpha^0_{mj}e^{\delta^\top Z(t)}$, now breathes in and out with
the covariate. So we keep a finite-dimensional, integrable system; we merely lose
the constant-coefficient (and hence transform-solvable) structure. ``Same state,
time-varying gain.''
\end{intuition}

The system \eqref{eq:ltv} is integrated by Runge--Kutta exactly like the
convolution case; the package provides it as \code{solve\_mbpp\_ltv}. Two checks
confirm correctness: with $Z\equiv 0$ (no modulation) it reproduces the LTI
multivariate solver to machine precision ($\sim10^{-16}$); and with a covariate
that, say, \emph{doubles} the excitation over a window, the intensity rises sharply
during the window and \emph{relaxes back} to the unmodulated stationary level
afterwards --- the signature of a stable forced system.

\begin{remark}[Power-law shapes: direct Volterra quadrature]
If $g$ is \emph{not} a finite sum of exponentials (e.g.\ power-law), the kernel is
\emph{not} separable and no finite ODE exists. One then solves
\eqref{eq:mbpp-nonconv} directly by discretising the Volterra integral on the
observation grid --- a triangular (lower-block-triangular) linear system solved by
forward substitution in $O(N^2 M^2)$ (the Nyström method,
Appendix~\ref{app:volterra}). One may also approximate $g$ by a sum of exponentials
(Prony) to restore a finite LTV ODE to any desired accuracy.
\end{remark}

\subsection{Fitting with excitation covariates}

The MBPP is \emph{still} a Poisson process with a deterministic (now
harder-to-compute) intensity, so the interval-censored likelihood~\eqref{eq:icll}
still applies verbatim: assemble $\Xi(o_{i-1},o_i]$ by integrating the LTV ODE
\eqref{eq:ltv} (or by quadrature), and minimise $\mathcal L_{\mathrm{IC}}$ over
$(\alpha^0,\theta,\delta,\gamma)$. Three things change relative to the convolution
case:
\begin{enumerate}[leftmargin=2.0em]
\item \textbf{More parameters.} Each excitation-covariate coefficient $\delta_{mj}$
adds dimensions; with $p$ covariates and an $M\times M$ network that is up to
$pM^2$ extra parameters. An $\ell_1$ penalty or low-rank/shared-$\delta$ structure
is advisable.
\item \textbf{Non-convexity.} The $\delta$ enter through $e^{\delta^\top Z}$
\emph{inside} the intensity, so the loss is non-convex (unlike the baseline-only
case, which is well-behaved). Use multiple restarts.
\item \textbf{Identification needs interaction.} A $\delta_{mj}$ is identified only
if component $j$ actually fires while $Z$ varies; without that co-occurrence the
excitation modulation leaves no fingerprint in the counts.
\end{enumerate}

\begin{intuition}
We have traded a closed-form, convex, transform-solvable problem for a
finite-dimensional but \emph{time-varying} and \emph{non-convex} one. The physics
is appealing --- ``contagion that is more contagious in some regimes than others''
--- and entirely doable; one simply pays for it in computation (integrate an ODE
instead of evaluating a formula) and in statistics (more parameters, a rougher
landscape, a stronger need for covariate variation). That trade is the precise
content of ``excitation covariates are harder.''
\end{intuition}

\subsection{Summary of the three regimes}

\begin{center}
\renewcommand{\arraystretch}{1.3}
\begin{tabular}{@{}lccc@{}}
\toprule
 & \S\ref{sec:no-covariates}/\ref{sec:baseline-cov} no/baseline cov.\ & \multicolumn{2}{c}{\S\ref{sec:excitation-cov} excitation covariates}\\
\cmidrule(l){3-4}
 & convolution & exp.\ shape & general shape\\
\midrule
MBPP equation & convolution Volterra & \multicolumn{2}{c}{general Volterra (non-convolution)}\\
Time-invariant? & yes (LTI) & no & no\\
Laplace/Fourier solution & yes & no & no\\
Reduces to finite ODE? & yes (constant coef.) & yes (\emph{time-varying} coef.) & no\\
How we compute $\xi,\Xi$ & closed form & integrate LTV ODE & Volterra quadrature\\
Fitting loss & convex(ish) & non-convex & non-convex\\
\bottomrule
\end{tabular}
\end{center}

\noindent
Reading the table left to right is reading this document: each added layer of
realism removes a layer of mathematical convenience, and the Volterra formulation
is what lets us track exactly which convenience is lost and what replaces it.


\appendix
\section{Volterra integral equations: a primer}
\label{app:volterra}

This appendix is self-contained and assumes no prior exposure to integral
equations. Everything the main text needs about Volterra equations is here.

\subsection{What is an integral equation, and why ``Volterra''?}

An \emph{integral equation} has the unknown function under an integral sign,
just as a \emph{differential} equation has it under a derivative. The kind we meet
is the \emph{linear Volterra equation of the second kind}:
\begin{equation}
\xi(t)=s(t)+\int_0^t K(t,u)\,\xi(u)\dd u,\qquad t\in[0,T].
\label{eq:volterra2}
\end{equation}
Here $s$ (the ``forcing'' or ``free term'') and $K$ (the ``kernel'') are given, and
$\xi$ is unknown. ``Second kind'' means the unknown also appears \emph{outside} the
integral (the lone $\xi(t)$ on the left); a \emph{first-kind} equation
$s(t)=\int_0^t K(t,u)\xi(u)\dd u$ has the unknown only inside and is much more
delicate (it is a deconvolution and is typically ill-posed).

The name \emph{Volterra} refers to the upper limit being the running time $t$:
the integral only looks at the \emph{past} $u\le t$. Contrast a \emph{Fredholm}
equation, whose integral runs over a fixed full interval $\int_0^T$. The Volterra
(causal) structure is what makes \eqref{eq:volterra2} so well-behaved: it is, in a
precise sense, ``lower-triangular in time.''

\begin{intuition}
Think of \eqref{eq:volterra2} as a feedback rule unrolling in time: the value
$\xi(t)$ ``now'' equals an external input $s(t)$ \emph{plus} a weighted echo of
everything $\xi(u)$ that happened before, with weight $K(t,u)$. Because the echo
only reaches back into the past, you can in principle sweep forward and build
$\xi$ left to right --- nothing in the future is needed to determine the present.
That causal, triangular character is the source of all the good news below.
\end{intuition}

\subsection{The operator form and the Neumann (resolvent) series}

Define the \emph{Volterra operator}
\begin{equation}
(\mathcal K\xi)(t):=\int_0^t K(t,u)\,\xi(u)\dd u,
\end{equation}
so \eqref{eq:volterra2} is $\xi=s+\mathcal K\xi$, i.e.\ $(\mathcal I-\mathcal
K)\xi=s$. Formally inverting the operator as a geometric series,
\begin{equation}
\xi=(\mathcal I-\mathcal K)^{-1}s=\sum_{n=0}^{\infty}\mathcal K^n s
= s+\mathcal K s+\mathcal K^2 s+\cdots
\label{eq:neumann}
\end{equation}
This is the \emph{Neumann series}. The operator $\mathcal K^n$ has its own kernel,
the \emph{iterated kernel} $K_n(t,u)$ (with $K_1=K$ and
$K_{n+1}(t,u)=\int_u^t K(t,r)K_n(r,u)\dd r$), and the sum
$R(t,u)=\sum_{n\ge1}K_n(t,u)$ is the \emph{resolvent kernel}. With it,
\begin{equation}
\xi(t)=s(t)+\int_0^t R(t,u)\,s(u)\dd u,
\label{eq:resolvent}
\end{equation}
so solving the equation amounts to computing $R$ \emph{once} and then a single
integral against any forcing $s$. $R$ satisfies the resolvent equations
$R=K+K\!\circ\!R=K+R\!\circ\!K$, where $\circ$ is the Volterra composition above.

\begin{intuition}
$\mathcal K s$ is the first-generation echo of the input, $\mathcal K^2 s$ the
echo of the echo, and so on; \eqref{eq:neumann} simply sums all generations. The
resolvent $R$ packages every generation into one kernel --- it is the
``total-feedback'' Green's function of the problem. (In the Hawkes setting $R$ is
literally the renewal density of the branching tree, \S\ref{sec:no-covariates}.)
\end{intuition}

\subsection{Existence and uniqueness: why Volterra is so forgiving}

\begin{theorem}[Existence and uniqueness on a finite horizon]\label{thm:vol-exist}
Suppose $K$ is bounded, $|K(t,u)|\le \bar K$ on the triangle $0\le u\le t\le T$,
and $s\in C[0,T]$. Then \eqref{eq:volterra2} has a unique continuous solution on
$[0,T]$, namely the (uniformly convergent) Neumann series \eqref{eq:neumann}.
\end{theorem}
\begin{proof}
By induction the iterated operator satisfies
$|(\mathcal K^n\xi)(t)|\le \bar K^{\,n}\,\|\xi\|_\infty\,t^n/n!$: indeed
$|(\mathcal K\xi)(t)|\le \bar K\|\xi\|_\infty t$, and
$|(\mathcal K^{n+1}\xi)(t)|\le \int_0^t \bar K\,|(\mathcal K^n\xi)(r)|\dd r
\le \bar K^{\,n+1}\|\xi\|_\infty\int_0^t r^n/n!\dd r=\bar K^{\,n+1}\|\xi\|_\infty t^{n+1}/(n+1)!$.
Hence $\|\mathcal K^n\|_{C[0,T]}\le(\bar K T)^n/n!\to 0$, the series
$\sum_n\mathcal K^n$ converges in operator norm, $(\mathcal I-\mathcal K)$ is
boundedly invertible, and the fixed point \eqref{eq:neumann} is unique.
\end{proof}

\begin{intuition}
The decisive estimate is the factorial $\bar K^n t^n/n!$. Each time the echo
passes through the kernel it must integrate over a \emph{shrinking} triangle of
``time still available,'' and $n$ such nested integrations contribute a $1/n!$ that
overwhelms any fixed gain $\bar K^n$. So on a finite horizon the loop is
\emph{always} invertible --- there is no smallness condition on the kernel at all.
This is the single most important structural fact, and it is why the MBPP equation
is solvable even when the excitation is covariate-modulated and time-varying
(\S\ref{sec:excitation-cov}): nothing in the proof used convolution structure.
\end{intuition}

\noindent
A smallness condition \emph{does} appear if one wants the resolvent globally on
$[0,\infty)$: for a convolution kernel $K(t,u)=\phi(t-u)$, Young's inequality gives
$\|\phi^{\conv n}\|_1\le\|\phi\|_1^n$, so $R=\sum_n\phi^{\conv n}\in L^1(0,\infty)$
iff $\|\phi\|_1=\kappa<1$ --- exactly subcriticality.

\subsection{Three solution strategies, by kernel type}

The structure of $K$ dictates the best method. The main text uses all three.

\paragraph{(i) Convolution kernel $K(t,u)=\phi(t-u)$ $\Rightarrow$ transforms / LTI.}
Then $\mathcal K\xi=\phi\conv\xi$ and the system is time-invariant. Laplace or
Fourier turn the convolution into a product (Appendix~\ref{app:transforms}):
$\hat\xi=\hat s/(1-\hat\phi)$. The resolvent is itself a convolution,
$R(t,u)=h(t-u)$ with $h=\sum_n\phi^{\conv n}$, and for exponential $\phi$ it is a
single exponential. \emph{This is \S\ref{sec:no-covariates}--\ref{sec:baseline-cov}.}

\paragraph{(ii) Separable (degenerate) kernel
$K(t,u)=\sum_{r=1}^{\rho}a_r(t)\,b_r(u)$ $\Rightarrow$ finite ODE.}
A finite-rank kernel collapses the integral equation onto $\rho$ ordinary
differential equations. Indeed, defining the states $y_r(t)=\int_0^t b_r(u)\xi(u)\dd u$,
the equation \eqref{eq:volterra2} becomes $\xi(t)=s(t)+\sum_r a_r(t)\,y_r(t)$, and
differentiating each state gives the closed ODE system
\begin{equation}
y_r'(t)=b_r(t)\,\xi(t)=b_r(t)\Bigl(s(t)+\textstyle\sum_{r'}a_{r'}(t)y_{r'}(t)\Bigr),
\qquad y_r(0)=0.
\label{eq:sep-ode}
\end{equation}
If the $a_r,b_r$ are constants (or pure exponentials), \eqref{eq:sep-ode} has
\emph{constant} coefficients --- a linear time-invariant ODE; if the $a_r$ carry an
extra $t$-dependence (as the covariate factor $e^{\delta^\top Z(t)}$ does), it is a
\emph{time-varying} linear ODE. \emph{An exponential triggering shape is separable
of rank one per component pair, which is precisely why
\S\ref{sec:excitation-cov} reduces to the LTV system \eqref{eq:ltv}.}

\paragraph{(iii) General kernel $\Rightarrow$ numerical quadrature (Nyström).}
When $K$ is neither convolution nor separable, discretise the integral on a grid
$0=u_0<\cdots<u_N=T$. Replacing $\int_0^{t_i}$ by a quadrature rule turns
\eqref{eq:volterra2} into a \emph{lower-triangular} linear system
\begin{equation}
\xi_i = s_i + \sum_{k\le i} w_{ik}\,K(t_i,u_k)\,\xi_k,\qquad i=0,\dots,N,
\end{equation}
which is solved by forward substitution in $O(N^2)$ (the diagonal term is handled
by isolating $\xi_i$). Convergence is at the order of the quadrature rule. This is
the fallback for, e.g., a power-law shape under covariate-modulated excitation.

\begin{intuition}
The three strategies are the same idea --- invert $\mathcal I-\mathcal K$ ---
specialised to how much structure $K$ has. Maximum structure (convolution): a
formula. Medium structure (finite rank): a handful of ODEs. No structure: march
forward in time filling in $\xi$ one step at a time, which the triangular
(causal) form always permits.
\end{intuition}

\begin{remark}[First-kind equations and why we avoid them]
The compensator's \emph{inverse} problem --- recovering $\phi$ from observed mean
behaviour --- is a first-kind (deconvolution) Volterra problem and is ill-posed
(small data errors blow up). The main text never inverts in this direction for
estimation; it always poses a well-posed second-kind forward solve inside a
likelihood, which is far more stable.
\end{remark}

\section{Laplace and Fourier methods for convolution kernels}
\label{app:transforms}

The transform methods of this appendix apply \emph{only} to the convolution case
$K(t,u)=\phi(t-u)$ (no/baseline covariates). Their failure for non-convolution
kernels is exactly what makes the excitation-covariate case
(\S\ref{sec:excitation-cov}) hard.

\subsection{The transforms and the convolution theorem}

The (one-sided) Laplace transform of $f$ is
$\Lap\{f\}(z)=\int_0^\infty e^{-zt}f(t)\dd t$, and the Fourier transform is
$\Four\{f\}(\omega)=\hat f(\omega)=\int_{-\infty}^\infty e^{-i\omega t}f(t)\dd t$.
The one property we need is that both turn convolution into multiplication:
\begin{equation}
\Lap\{\phi\conv\xi\}=\Lap\{\phi\}\,\Lap\{\xi\},\qquad
\widehat{\phi\conv\xi}=\hat\phi\,\hat\xi .
\end{equation}

\begin{intuition}
A convolution mixes every past value of $\xi$ into the present through a
\emph{shift-invariant} weight. The transforms re-express $\xi$ in terms of pure
oscillations $e^{i\omega t}$ (or modes $e^{zt}$), and a shift-invariant system acts
on each such mode merely by scaling it. So in the transform domain the tangled
convolution becomes an independent, per-frequency multiplication --- the system is
``diagonalised.''
\end{intuition}

\subsection{Solving the convolution Volterra equation}

Apply $\Lap$ to $\xi=s+\phi\conv\xi$:
$\Lap\{\xi\}=\Lap\{s\}+\Lap\{\phi\}\Lap\{\xi\}$, hence
\begin{equation}
\Lap\{\xi\}(z)=\frac{\Lap\{s\}(z)}{1-\Lap\{\phi\}(z)},
\qquad
\hat\xi(\omega)=R(\omega)\,\hat s(\omega),\quad R(\omega)=\frac{1}{1-\hat\phi(\omega)} .
\label{eq:transfer}
\end{equation}
The multiplier $R$ is the \emph{transfer function}; inverting the transform
recovers $\xi$. $R=1/(1-\hat\phi)$ is the frequency-domain face of the geometric
series $\sum_n\hat\phi^{\,n}$ --- the resolvent of Appendix~\ref{app:volterra}.

\begin{example}[Exponential kernel]
For $\phi(t)=\kappa\theta e^{-\theta t}$, $\Lap\{\phi\}(z)=\kappa\theta/(z+\theta)$,
so
\begin{equation*}
1-\Lap\{\phi\}=\frac{z+(1-\kappa)\theta}{z+\theta},\qquad
R(\omega)=\frac{\theta+i\omega}{(1-\kappa)\theta+i\omega}.
\end{equation*}
This is a one-pole filter; its pole at $-(1-\kappa)\theta$ is precisely the rate of
the ODE in Appendix~\ref{app:ode}, and inverting $\Lap\{\phi\}/(1-\Lap\{\phi\})$
gives the single-exponential resolvent $h(t)=\kappa\theta e^{(\kappa-1)\theta t}$.
Time domain, frequency domain, and state space are three faces of one operator.
\end{example}

In the multivariate case the same computation gives
$\hat\xi(\omega)=(I-\hat\Phi(\omega))^{-1}\hat s(\omega)$, and at $\omega=0$ the
matrix $\hat\Phi(0)=G$ is the branching matrix, recovering the stationary law
$\bar\xi=(I-G)^{-1}\bar\mu$ (Proposition~\ref{prop:stationary}).

\subsection{The spectral (neural-operator) viewpoint, and why it is exact here}

Because \eqref{eq:transfer} is a pointwise multiply in frequency, the solution
operator $s\mapsto\xi$ is \emph{diagonal} under the Fourier transform. A single
linear Fourier-neural-operator layer computes exactly
$\xi=\Four^{-1}\{R\cdot\Four\{s\}\}$, so with $R=1/(1-\hat\phi)$ it represents the
MBPP operator with \emph{no} approximation; conversely, \emph{learning} $R(\omega)$
from input/output pairs is nonparametric identification of the operator (and hence
of $\phi$ via $\hat\phi=1-1/R$).

\subsection{Why transforms fail under excitation covariates}

All of the above hinges on $K(t,u)=\phi(t-u)$. When the excitation is
covariate-modulated, $K(t,u)=\alpha^0 e^{\delta^\top Z(t)}g(t-u)$ depends on $t$
and $u$ \emph{separately}; the convolution theorem no longer applies, there is no
single transfer function, and the operator is \emph{not} diagonalised by Fourier.
One must return to the time domain --- the LTV ODE or Volterra quadrature of
Appendix~\ref{app:volterra}. This is the precise mathematical price of
\S\ref{sec:excitation-cov}.

\begin{intuition}
Transforms diagonalise \emph{shift-invariant} systems. A covariate that changes the
gain over time breaks shift-invariance --- the system treats ``now'' differently
from ``then'' --- so no single set of frequency weights can describe it. The cure is
to stop trying to diagonalise and instead integrate the dynamics forward in time.
\end{intuition}

\section{ODE and state-space methods}
\label{app:ode}

The bridge between integral equations and differential equations: when the kernel
is built from exponentials, the Volterra equation \emph{is} an ODE in disguise.

\subsection{The exponential kernel gives a first-order ODE}

Take the convolution MBPP equation $\xi=s+\phi\conv\xi$ with
$\phi(t)=\kappa\theta e^{-\theta t}$, and split off the endogenous part
$y:=\xi-s=\phi\conv\xi$. Differentiating
$y(t)=\kappa\theta\int_0^t e^{-\theta(t-u)}\xi(u)\dd u$ (Leibniz rule),
\begin{equation}
y'(t)=\kappa\theta\,\xi(t)-\theta\,y(t)=(\kappa-1)\theta\,y(t)+\kappa\theta\,s(t),
\qquad y(0)=0,\quad \xi=s+y .
\label{eq:exp-ode}
\end{equation}
So $\xi$ obeys a \emph{first-order linear ODE forced by the baseline $s$}, with a
\emph{stable} pole at $(\kappa-1)\theta<0$. Solving \eqref{eq:exp-ode} with the
integrating factor reproduces the single-exponential resolvent
$h(t)=\kappa\theta e^{(\kappa-1)\theta t}$.

\begin{intuition}
A single memory timescale is a single pole, i.e.\ a first-order filter, i.e.\ a
first-order ODE. One number $\kappa$ sets the feedback gain and one number $\theta$
sets the relaxation rate $(1-\kappa)\theta$: push the system and it returns to
baseline at that rate.
\end{intuition}

\subsection{Rational Laplace transform $\Leftrightarrow$ finite ODE}

\begin{theorem}[Realisation]\label{thm:rational}
The convolution Volterra equation $\xi=s+\phi\conv\xi$ is equivalent to a
finite-order linear constant-coefficient ODE (forced by $s$ and its derivatives)
\emph{if and only if} $\phi$ has a strictly proper rational Laplace transform ---
equivalently, iff $\phi$ is a finite exponential-polynomial
$\sum_q p_q(t)e^{-b_q t}$.
\end{theorem}
\begin{proof}
($\Leftarrow$) If $\Lap\{\phi\}=P/Q$ with $\deg P<\deg Q=:d$, then from
\eqref{eq:transfer} $\Lap\{\xi\}=\frac{Q}{Q-P}\Lap\{s\}$, so
$(Q-P)(z)\Lap\{\xi\}=Q(z)\Lap\{s\}$. As $Q-P$ has degree $d$, reading
$z^k\leftrightarrow d^k/\dd t^k$ yields the order-$d$ ODE
$(Q-P)(\tfrac{d}{dt})\xi=Q(\tfrac{d}{dt})s$. ($\Rightarrow$) Conversely a finite ODE
makes the transfer function $Q/(Q-P)$ rational, hence
$\Lap\{\phi\}=1-(Q-P)/Q=P/Q$ rational; and a function has rational Laplace transform
iff it is an exponential-polynomial (partial fractions and inverse Laplace).
\end{proof}

\begin{intuition}
A memory built from finitely many decaying modes can be tracked by finitely many
state variables --- ``rational transfer function $\Leftrightarrow$
finite-dimensional machine'' (this is realisation theory). A power-law memory is
\emph{not} rational: it carries infinitely many timescales, so no finite ODE exists
(one then uses Volterra quadrature, Appendix~\ref{app:volterra}).
\end{intuition}

\subsection{Sums of exponentials: the state-space realisation}

For $\phi(\tau)=\sum_q a_q e^{-b_q\tau}$ the natural states are
$y_q=\bigl(a_q e^{-b_q\cdot}\bigr)\conv\xi$, giving the constant-coefficient system
\begin{equation}
y_q'=a_q\,\xi-b_q\,y_q,\qquad \xi=s+\sum_q y_q,
\end{equation}
i.e.\ $\dot u=Au+a\,s$, $\xi=s+\mathbf 1^\top u$ with $A=a\mathbf 1^\top-\diag(b)$.
The multivariate version (kernel matrix) carries an index pair $(m,j)$ and gives the
$M^2$-state system \eqref{eq:statespace}, implemented as
\code{solve\_mbpp\_ode\_multivariate}. Since sums of exponentials are dense, this
realises (or approximates) essentially any kernel.

\subsection{Numerical integration and the time-varying case}

These ODEs are integrated by a standard explicit Runge--Kutta (RK4) sweep on the
observation grid; the compensator $\Xi=\int_0^t\xi$ is accumulated alongside (an
extra trapezoidal state). The constant-coefficient systems are autonomous in their
coefficients; the \emph{covariate-modulated} system \eqref{eq:ltv} is identical in
form but with the readout gain $\alpha^0_{mj}e^{\delta_{mj}^\top Z(t)}$ depending on
$t$ --- a \emph{linear time-varying} (LTV) ODE.

\begin{proposition}[LTV well-posedness]
If the coefficient maps $t\mapsto A(t)$ and $t\mapsto a(t)$ are continuous (here
through the continuous covariate path $Z(t)$), the LTV system $\dot u=A(t)u+a(t)s(t)$
has a unique solution on $[0,T]$ for every initial condition, given by the
time-ordered (Peano--Baker) series; numerically, RK4 converges at its usual order.
\end{proposition}

\begin{intuition}
A time-varying linear ODE is no harder to \emph{integrate} than a constant one ---
you just re-read the coefficients at each step. What you lose is the \emph{closed
form}: there is no single matrix exponential $e^{At}$ when $A$ depends on $t$ (the
$A(t)$ at different times need not commute), which is the differential-equations
echo of ``no constant transfer function'' from Appendix~\ref{app:transforms}. So
\code{solve\_mbpp\_ltv} integrates \eqref{eq:ltv} step by step, and that is exactly
how the excitation-covariate model of \S\ref{sec:excitation-cov} is computed.
\end{intuition}

\section{Estimation, Bregman geometry, and identifiability}
\label{app:estimation}

\subsection{IC-LL is a Bregman divergence}

Write the interval counts $\sfont C=(\sfont C_i)$ and compensators
$\Xi(\theta)=(\Xi_i(\theta))$. A \emph{Bregman divergence} for a strictly convex
generator $\varphi$ is
$B_\varphi(x,y)=\varphi(x)-\varphi(y)-\langle x-y,\nabla\varphi(y)\rangle\ge0$.

\begin{proposition}\label{prop:kl}
With the generalized-KL generator $\varphi(x)=\sum_i(x_i\log x_i-x_i)$,
$\mathcal L_{\mathrm{IC}}(\theta)=\KL(\sfont C,\Xi(\theta))+\Gamma(\sfont C)$ with
$\Gamma$ free of $\theta$; hence
$\argmin_\theta\mathcal L_{\mathrm{IC}}=\argmin_\theta\KL(\sfont C,\Xi(\theta))$.
The squared-error generator $\tfrac12\|x\|^2$ instead yields the least-squares loss
$\sum_i(\sfont C_i-\Xi_i)^2$.
\end{proposition}

\begin{theorem}[Bregman $\leftrightarrow$ exponential family; Banerjee et al.]
For a regular exponential family with cumulant $\psi$ and mean $\mu=\nabla\psi$, the
negative log-likelihood equals $B_{\psi^\star}(x,\mu)$ up to a constant, where
$\psi^\star$ is the convex conjugate. Poisson gives $\psi^\star(\mu)=\mu\log\mu-\mu$
(so $B=\KL$); Gaussian gives $\psi^\star(\mu)=\mu^2/2$ (so $B=$ squared error).
\end{theorem}

\begin{intuition}
Maximum likelihood in any exponential family is ``move the model mean as close as
possible to the data, in the geometry set by the family's variance function.''
Choosing the divergence \emph{is} choosing the noise model: $x\log x\to$ Poisson
$\to$ KL $\to$ IC-LL; $\tfrac12x^2\to$ Gaussian $\to$ least squares. Because a
Poisson bucket has variance equal to its mean, KL correctly down-weights
large-count buckets by $1/\Xi_i$, whereas least squares treats all buckets alike
and is dominated by the loud ones --- so KL is the more robust loss when counts span
orders of magnitude (bursts, near criticality).
\end{intuition}

\subsection{Consistency, information, and goodness of fit}

\begin{theorem}[Asymptotics of the MLE; Ogata]
For a stationary, ergodic, subcritical process with true $\theta_0$ in the interior
of a compact, identifiable $\Theta$, twice-differentiable log-intensity with
integrable dominating envelopes, and finite nonsingular per-time information
$\mathfrak i(\theta_0)$, the MLE satisfies $\hat\theta_T\to\theta_0$ a.s.\ and
$\sqrt T(\hat\theta_T-\theta_0)\to\mathcal N(0,\mathfrak i(\theta_0)^{-1})$.
\end{theorem}
The information is
$\mathcal I(\theta)=\E\!\int_0^T\frac{\partial_\theta\lambda\,\partial_\theta\lambda^\top}{\lambda}\dd t$
(event-time) or
$\mathcal I^{\mathrm{IC}}_{jk}=\sum_i\partial_{\theta_j}\Xi_i\,\partial_{\theta_k}\Xi_i/\Xi_i$
(interval-censored); standard errors come from inverting the observed information
(the Hessian of $-\ell$), and Cram\'er--Rao says no unbiased estimator beats
$\mathcal I^{-1}$.

\begin{theorem}[Time-rescaling; goodness of fit]
If $\Lambda$ is continuous, strictly increasing with $\Lambda(\infty)=\infty$, the
rescaled times $\Lambda(t_k)$ form a unit-rate Poisson process; equivalently the
rescaled gaps $\Lambda(t_k)-\Lambda(t_{k-1})$ are i.i.d.\ $\mathrm{Exp}(1)$.
\end{theorem}

\begin{intuition}
Measure time not in seconds but in \emph{accumulated compensator} --- ``operational
time'' that ticks fast when the intensity is high. In that clock the events look
like a plain unit-rate Poisson process, giving a model-free residual check
(rescaled gaps should be $\mathrm{Exp}(1)$).
\end{intuition}

\subsection{The $\kappa$--$\theta$ ridge, and the cost of covariates}

\begin{proposition}[Weak identifiability of the timescale]
In $\phi=\alpha e^{-\beta t}$, $\alpha=\kappa\theta$, $\beta=\theta$, the
$(\alpha,\beta)$ information matrix is typically ill-conditioned, with its
small-eigenvalue direction aligned with curves of constant $\alpha/\beta=\kappa$.
Hence $\kappa$ is well identified while $\theta$ is weak, and the more so the
coarser the censoring: as the bucket width grows the timing-borne sensitivity
$\partial_\theta\Xi_i\to0$ while $\partial_\kappa\Xi_i$ stays order one.
\end{proposition}

\begin{intuition}
The counts pin down ``how much'' (the gain $\kappa$) but not ``how fast'' (the
timescale $\theta$), because a bucket integrates away the within-bucket timing that
carries $\theta$. So treat $\kappa$ and the covariate effects $\gamma$ as the
trustworthy estimands; fix or prior-constrain $\theta$ from domain knowledge and
report it with wide intervals.
\end{intuition}

\paragraph{Excitation covariates add non-convexity.} Baseline covariates
(\S\ref{sec:baseline-cov}) enter the compensator linearly through a log-link in the
\emph{forcing} and keep the problem well-behaved; the covariate effect $\gamma$ is
recovered tightly. Excitation covariates (\S\ref{sec:excitation-cov}) enter through
$e^{\delta^\top Z}$ \emph{inside} the intensity dynamics, so the likelihood is
non-convex in $\delta$, with up to $pM^2$ extra parameters and a real need for
covariate--event co-occurrence to identify each $\delta_{mj}$. This is the
statistical half of ``excitation covariates are harder''; the computational half is
that $\Xi$ now comes from integrating an LTV ODE rather than from a formula.


\begin{thebibliography}{9}
\small
\bibitem{Hawkes1971} A.~G.~Hawkes (1971). Spectra of some self-exciting and mutually exciting point processes. \emph{Biometrika} 58(1).
\bibitem{HawkesOakes1974} A.~G.~Hawkes, D.~Oakes (1974). A cluster process representation of a self-exciting process. \emph{J.~Appl.~Probab.} 11(3).
\bibitem{DassiosZhao2013} A.~Dassios, H.~Zhao (2013). Exact simulation of Hawkes process with exponentially decaying intensity. \emph{Electron.~Commun.~Probab.} 18.
\bibitem{Rizoiu2017} M.-A.~Rizoiu et al. (2017). Expecting to be HIP: Hawkes intensity processes for social media popularity. \emph{WWW}.
\bibitem{Banerjee2005} A.~Banerjee, S.~Merugu, I.~Dhillon, J.~Ghosh (2005). Clustering with Bregman divergences. \emph{JMLR} 6.
\bibitem{Ogata1978} Y.~Ogata (1978). The asymptotic behaviour of maximum likelihood estimators for stationary point processes. \emph{Ann.~Inst.~Statist.~Math.} 30.
\bibitem{ICHawkes2022} M.-A.~Rizoiu, A.~Soen, S.~Li, P.~Calderon, L.~J.~Dong, A.~K.~Menon, L.~Xie (2022). Interval-censored Hawkes processes. \emph{JMLR} 23(1):1--84.
\bibitem{GLS1990} G.~Gripenberg, S.-O.~Londen, O.~Staffans (1990). \emph{Volterra Integral and Functional Equations}. Cambridge Univ.~Press.
\bibitem{ChenChen1995} T.~Chen, H.~Chen (1995). Universal approximation to nonlinear operators by neural networks. \emph{IEEE Trans.~Neural Networks} 6(4).
\end{thebibliography}

\end{document}
